\section*{Hoofdstuk \ref{ch:b2:genome-diversification} --- Mutaties en diversificatie van het genoom}

\begin{solution}{exo:b2:genome-diversification:1}
GAG: stil (nog steeds glutamaat), een transitie (A$\to$G). GTA: missense
(valine), een transversie (A$\to$T). TAA: nonsense (stop), een
transversie (G$\to$T). Invoegen van een base: een \omterm{def:b2:genome-diversification:mutation}{leesraamverschuiving},
die elk volgend codon herschrijft.
\end{solution}

\begin{solution}{exo:b2:genome-diversification:2}
$U = 5\times 10^{-10}\times 4.6\times 10^{6} = 2.3\times 10^{-3}$: één
deling op ongeveer 430 levert een \omterm{def:b2:genome-diversification:mutation}{mutatie} op. In $10^{9}$ cellen:
$2.3\times 10^{6}$ nieuwe \omterm{def:b2:genome-diversification:mutation}{mutaties} per generatie.
\end{solution}

\begin{solution}{exo:b2:genome-diversification:3}
\omterm{def:b2:genome-diversification:hgt}{Transformatie}: opname van vrij DNA uit gelyseerde cellen, vereist een
competente ontvanger; draagt chromosoomfragmenten over (kapselgenen bij
pneumokokken). Conjugatie: contact van cel tot cel via een pilus, vereist
een conjugatief \omterm{def:b2:unicellular-diversity:prokaryote}{plasmide} in de donor; draagt \omterm{def:b2:unicellular-diversity:prokaryote}{plasmiden} over
(resistentiegenen) en, vanuit een Hfr-stam, chromosomale genen.
\omterm{def:b2:genome-diversification:hgt}{Transductie}: een bacteriofaag die gastheer-DNA heeft verpakt; draagt een
willekeurig fragment over (gegeneraliseerd) of de genen naast de plaats
van de profaag (gespecialiseerd).
\end{solution}

\begin{solution}{exo:b2:genome-diversification:4}
\omterm{def:b2:genome-diversification:polyploidy}{Autopolyploïd}: extra sets van één soort, via een diploïde gameet
(tetraploïde aardappel, veel autotetraploïde klavers). \omterm{def:b2:genome-diversification:polyploidy}{Allopolyploïd}: de
sets van twee soorten, gecombineerd door hybridisatie en daarna
verdubbeld (broodtarwe, katoen, tabak). De diploïde hybride AB heeft geen
twee gelijke chromosomen, dus kan geen enkel chromosoom bij meiose I
paren; de chromosomen segregeren willekeurig en de gameten zijn
ongebalanceerd en niet levensvatbaar.
\end{solution}

\begin{solution}{exo:b2:genome-diversification:5}
$2\times 3.2\times 10^{9}\times 1.2\times 10^{-8} = 77$ nieuwe \omterm{def:b2:genome-diversification:mutation}{mutaties}.
In coderende sequentie: $77\times 0.015 = 1.2$; met verandering van een
aminozuur: $1.2\times 0.7 = 0.8$ --- de meeste kinderen dragen ongeveer
één nieuwe aminozuursubstitutie.
\end{solution}

\begin{solution}{exo:b2:genome-diversification:6}
Gebeurtenissen: $m(N - N_0) = 2\times 10^{-8}\times 10^{9} = 20$. Mutante
cellen: $mN\ln(N/N_0) = 20\times \ln 10^{6} = 20\times 13.8 = 276$.
$p_0 = e^{-20} \approx 2\times 10^{-9}$: elke cultuur heeft mutanten,
dus het tellen van culturen zonder mutant meet niets. Men heeft $m(N -
N_0) \approx 1$ nodig, d.w.z.\ $N \approx 5\times 10^{7}$, waarbij $p_0 \approx
0.37$ en enkele tientallen culturen $m$ tot op een factor twee geven.
\end{solution}

\begin{solution}{exo:b2:genome-diversification:7}
Uracil hoort niet in DNA thuis, dus kan een uracil-DNA-glycosylase elk
uracil dat het vindt zonder dubbelzinnigheid verwijderen. Thymine is een
normale base: nadat 5-methylcytosine is gedeamineerd, wordt de
G$\cdot$T-mismatch wel gezien, maar de herstelmachinerie kan niet weten
welke streng fout is en ``corrigeert'' de helft van de tijd de G,
waardoor een C$\to$T-transitie wordt vastgelegd. Gemethyleerde
CpG-plaatsen muteren daardoor ongeveer tien keer sneller dan andere
plaatsen, en in de loop van de evolutie zijn ze omgezet in TpG en CpA:
genomen van zoogdieren bevatten slechts ongeveer een vijfde van het CpG
dat de basensamenstelling voorspelt.
\end{solution}

\begin{solution}{exo:b2:genome-diversification:8}
\qty{46}{kb} per minuut. Posities: 368, 782, 1150 en
\qty{1840}{kb}; afstanden 414, 368 en \qty{690}{kb}.
\end{solution}

\begin{solution}{exo:b2:genome-diversification:9}
De tweede kopie van het ene chromosoom paart met de eerste kopie van het
andere; een crossing-over tussen beide levert één chromatide met
kopie~1, kopie~2 en kopie~2$'$ (drie genen) en één met alleen
kopie~1$'$. Iemand met van elke ouder een chromosoom met één gen heeft
twee $\alpha$-genen in plaats van vier ($-\alpha/-\alpha$):
$\alpha$-ketens worden op halve snelheid gemaakt, de rode bloedcellen zijn
klein en bleek, met milde of geen bloedarmoede --- dragerschap van
$\alpha$-thalassemie.
\end{solution}

\begin{solution}{exo:b2:genome-diversification:10}
Gameten: $n = 21$ (één set elk van A, B, D). Hybriden: AB, 14
chromosomen; ABD, 21. Broodtarwe $\times$ emmertarwe: AABBD, 35
chromosomen --- de sets A en B paren, de set D heeft geen partners,
en de plant is grotendeels steriel. Elke verdubbeling gaf elk chromosoom
een identieke homoloog, zodat bivalenten ontstaan en de gameten
gebalanceerd zijn; terugkruisen met een ouder levert planten met één
ongepaarde set, waarvan de gameten ongebalanceerd zijn, zodat de
hexaploïde reproductief geïsoleerd is van haar voorouders --- een soort
in één stap.
\end{solution}

\begin{solution}{exo:b2:genome-diversification:11}
Logistische groei: $R(t) = N/\bigl(1 + (N - 1)e^{-\gamma Nt}\bigr)$.
De helft van de populatie wanneer $(N - 1)e^{-\gamma Nt} = 1$: $t = \ln(N -
1)/\gamma N = 20.7/0.5 = \qty{41}{h}$. Zonder het antibioticum groeien
plasmidevrije cellen \qty{5}{\%} sneller, zodat de verhouding van dragers
tot niet-dragers afneemt als $e^{-st}$ met $s = 0.05\ln 2 =
\qty{0.035}{h^{-1}}$ (\omterm{def:b2:unicellular-diversity:division}{generatietijd} \qty{1}{h}): van
\qty{99}{\%} naar \qty{1}{\%} duurt $\ln(99^2)/0.035 = \qty{260}{h}$,
elf dagen --- en voortdurende overdracht kan het \omterm{def:b2:unicellular-diversity:prokaryote}{plasmide} onbeperkt op een
lage frequentie houden, klaar voor de volgende kuur met het middel.
\end{solution}

\begin{solution}{exo:b2:genome-diversification:12}
De \omterm{thm:b2:genome-diversification:luria}{fluctuatietest} toonde aan dat resistentiemutaties vóór en zonder het
selecterende middel ontstaan, op willekeurige momenten, zodat \omterm{def:b2:genome-diversification:mutation}{mutatie}
blind is voor wat nuttig zou zijn; het replicaplaten toonde hetzelfde
zonder de geselecteerde cellen ook maar bloot te stellen. Evolutie is
niet willekeurig, omdat selectie deze blinde varianten consequent en
cumulatief naar hun gevolgen sorteert. De enige plaats waar het toeval
zelf wordt gevormd, is de snelheid: lijnen met te veel of te weinig
fouten verliezen, zodat de \omterm{thm:b2:genome-diversification:rate}{mutatiesnelheid} een aangepast compromis is
--- maar wat elke \omterm{def:b2:genome-diversification:mutation}{mutatie} doet, blijft onvoorspelbaar en ongekozen.
\end{solution}

\begin{solution}{pb:b2:genome-diversification:1}
\textbf{1.} Som 119, gemiddelde $5.95$; gemiddelde van de kwadraten
$11489/20 = 574$; variantie $574 - 35 = 540$.
\textbf{2.} Som 324, gemiddelde $16.2$; gemiddelde van de kwadraten
$5462/20 = 273$; variantie $273 - 262 = 11$.
\textbf{3.} De monsters uit één cultuur zijn poissonverdeeld (variantie
dicht bij het gemiddelde): het zijn willekeurige trekkingen uit één vaste
mutantfractie. De onafhankelijke culturen hebben een variantie die negentig
keer het gemiddelde is: hun mutanten ontstonden op verschillende,
willekeurige momenten tijdens de groei, en hoe vroeger de \omterm{def:b2:genome-diversification:mutation}{mutatie}, hoe
groter haar kloon.
\textbf{4.} Een jackpot: een \omterm{def:b2:genome-diversification:mutation}{mutatie} in een van de eerste delingen,
waarvan de kloon daarna met de cultuur meegroeide tot honderd cellen.
\textbf{5.} Veertien van de twintig: $p_0 = 0.70$; $m(N - N_0) = -\ln 0.70
= 0.36$; $m = 0.36/2\times 10^{8} = \qty{1.8e-9}{}$ per deling.
\textbf{6.} $1.8\times 10^{-9}\times 2\times 10^{8}\times \ln(2\times
10^{6}) = 0.36\times 14.5 = 5.2$ mutante cellen; waargenomen $5.95$ ---
overeenstemming binnen de ruis van twintig culturen.
\textbf{7.} Geen enkel monster was leeg (alle delen de mutanten van de
moedercultuur), dus geeft $p_0 = 0$ geen schatting; de monsters meten de
mutantfractie van één cultuur, die afhangt van het toevallige moment
waarop haar \omterm{def:b2:genome-diversification:mutation}{mutaties} optraden.
\textbf{8.} $u = m/20 = \qty{9e-11}{}$ per basenpaar per replicatie.
\textbf{9.} $U = 9\times 10^{-11}\times 4.6\times 10^{6} =
\qty{4.1e-4}{}$.
\textbf{10.} $10^{12}$ cellen $\times$ $4.1\times 10^{-4}$ $= 4.1\times
10^{8}$ nieuwe \omterm{def:b2:genome-diversification:mutation}{mutaties}.
\textbf{11.} $3\times 4.6\times 10^{6} = 1.4\times 10^{7}$ mogelijke
veranderingen; elk ontstaat $4.1\times 10^{8}/1.4\times 10^{7} \approx 30$
keer.
\textbf{12.} Elke verandering van één base, ook elke verandering die
resistentie geeft tegen een middel dat door één \omterm{def:b2:genome-diversification:mutation}{mutatie} wordt verslagen,
is al in tientallen cellen aanwezig voordat het middel wordt toegevoegd;
het middel selecteert ze alleen.
\textbf{13.} Twee onafhankelijke \omterm{def:b2:genome-diversification:mutation}{mutaties} in dezelfde cel treden op met
$m^2 \approx (2\times 10^{-9})^2 = 4\times 10^{-18}$ per deling: in
$10^{12}$ cellen $4\times 10^{-6}$ per generatie --- een dubbel
resistente cel bestaat vrijwel nooit vooraf, en daarom wordt tuberculose
met meerdere middelen tegelijk behandeld.
\textbf{14.} Met $a = N - 1$ en $E = e^{-\gamma Nt}$: $R = N/(1 +
aE)$, $\mathrm{d}R/\mathrm{d}t = \gamma N\cdot NaE/(1 + aE)^2$; en
$\gamma R(N - R) = \gamma\,\frac{N}{1 + aE}\cdot\frac{NaE}{1 + aE}$,
hetzelfde; $R(0) = N/(1 + N - 1) = 1$.
\textbf{15.} $(N - 1)E = 1$: $t = \ln(N - 1)/\gamma N = 20.7/0.5 =
\qty{41}{h}$.
\textbf{16.} $1 + aE = 1/0.99$, $aE = 0.0101$: $t = \ln(10^{9}/0.0101)
/0.5 = (20.7 + 4.6)/0.5 = \qty{51}{h}$.
\textbf{17.} $\gamma (N/2)^2 = \gamma N\cdot N/4 = 0.5\times 2.5\times
10^{8} = 1.25\times 10^{8}$ overdrachten per uur.
\textbf{18.} Groeisnelheden $\ln 2 = \qty{0.693}{h^{-1}}$ en
\qty{0.762}{h^{-1}}; de verhouding mutant/gevoelig groeit als $e^{st}$
met $s = \qty{0.069}{h^{-1}}$; van $1/N$ naar 1 duurt $\ln N/s =
20.7/0.069 = \qty{300}{h}$, twaalf dagen.
\textbf{19.} Overdracht is evenredig met het product van dragers en
niet-dragers en verloopt in het tempo van de ontmoetingen, niet van de
delingen; afstamming wordt beperkt door het groeivoordeel van de mutant,
dat klein is. Een \omterm{def:b2:unicellular-diversity:prokaryote}{plasmide} doorkruist de populatie in twee dagen, een
bevoordeelde mutant in twee weken.
\textbf{20.} $U = 4.1\times 10^{-2}$; schadelijk: $0.041\times 0.88\times
0.3 = 0.011$ per generatie.
\textbf{21.} Gemiddelde $2.2$: $e^{-2.2} = 0.11$, één nakomeling op negen
ongeschonden.
\textbf{22.} Wildtype: $1.1\times 10^{-4}$ per generatie, $0.022$ in
200: $e^{-0.022} = 0.98$.
\textbf{23.} Gemiddeld $520$ mutante cellen; $p_0 = e^{-36} \approx
10^{-16}$: geen enkele lege cultuur.
\textbf{24.} Onder intense, wisselende selectie (de roterende antibiotica
van een ziekenhuis) weegt het honderdvoudige aanbod aan nieuwe varianten
op tegen de last, en liften mutators mee met de resistenties die ze
voortbrengen. In een stabiele omgeving valt er niets te winnen en per
honderd generaties één schadelijke \omterm{def:b2:genome-diversification:mutation}{mutatie} te verliezen, zodat
mutatorlijnen worden weggeselecteerd.
\textbf{25.} $m = \qty{1.8e-9}{}$ per deling naar resistentie, $u =
\qty{9e-11}{}$ per basenpaar per replicatie, $U = \qty{4.1e-4}{}$
\omterm{def:b2:genome-diversification:mutation}{mutaties} per genoom per replicatie; het \omterm{def:b2:unicellular-diversity:prokaryote}{plasmide} bereikt de helft van de
darm in \qty{41}{h} en \qty{99}{\%} ervan in \qty{51}{h}.
\end{solution}
